\documentclass[11pt,a4paper]{article}

% ── Packages ─────────────────────────────────────────────────────────────────
\usepackage[a4paper, top=22mm, bottom=22mm, left=22mm, right=22mm]{geometry}
\usepackage[T1]{fontenc}
\usepackage{lmodern}
\usepackage{microtype}
\usepackage[table]{xcolor}
\usepackage{booktabs}
\usepackage{tabularx}
\usepackage{array}
\usepackage{parskip}
\usepackage{enumitem}
\usepackage{titlesec}
\usepackage{mdframed}
\usepackage{multicol}
\usepackage{fancyhdr}
\usepackage{graphicx}
\usepackage{hyperref}

% ── Colours ───────────────────────────────────────────────────────────────────
\definecolor{darkgreen}{HTML}{1B5E20}
\definecolor{midgreen}{HTML}{2E7D32}
\definecolor{lightgreen}{HTML}{43A047}
\definecolor{palegreen}{HTML}{F1F8E9}
\definecolor{teal}{HTML}{00696E}
\definecolor{amber}{HTML}{E65C00}
\definecolor{crimson}{HTML}{C62828}
\definecolor{slate}{HTML}{37474F}
\definecolor{ink}{HTML}{212121}
\definecolor{muted}{HTML}{78909C}

% ── Typography ────────────────────────────────────────────────────────────────
\setlist[itemize]{leftmargin=1.4em, itemsep=1pt, parsep=0pt, topsep=3pt}
\setlist[enumerate]{leftmargin=1.6em, itemsep=1pt, parsep=0pt, topsep=3pt}

\titleformat{\section}{\large\bfseries\color{darkgreen}}{}{0em}{}[\color{lightgreen}\rule{\linewidth}{0.6pt}]
\titleformat{\subsection}{\normalsize\bfseries\color{slate}}{}{0em}{}
\titlespacing{\section}{0pt}{14pt}{5pt}
\titlespacing{\subsection}{0pt}{8pt}{2pt}

% ── Header / Footer ───────────────────────────────────────────────────────────
\pagestyle{fancy}
\fancyhf{}
\renewcommand{\headrulewidth}{0pt}
\fancyhead[L]{\small\color{muted}300,000 Streets of Melbourne --- Technical Overview}
\fancyhead[R]{\small\color{muted}Regen Melbourne \texttimes{} RMIT University}
\fancyfoot[C]{\small\color{muted}\thepage}

% ── Custom environments ───────────────────────────────────────────────────────
\newmdenv[
  backgroundcolor=palegreen,
  linecolor=midgreen,
  linewidth=1.4pt,
  leftline=true, rightline=false, topline=false, bottomline=false,
  innerleftmargin=10pt, innerrightmargin=10pt,
  innertopmargin=7pt, innerbottommargin=7pt,
  skipabove=8pt, skipbelow=6pt
]{insightbox}

\newmdenv[
  backgroundcolor=crimson!8,
  linecolor=crimson,
  linewidth=1.4pt,
  leftline=true, rightline=false, topline=false, bottomline=false,
  innerleftmargin=10pt, innerrightmargin=10pt,
  innertopmargin=7pt, innerbottommargin=7pt,
  skipabove=8pt, skipbelow=6pt
]{findingbox}

% ── Column types ──────────────────────────────────────────────────────────────
\newcolumntype{L}[1]{>{\raggedright\arraybackslash}p{#1}}
\newcolumntype{C}[1]{>{\centering\arraybackslash}p{#1}}

% =============================================================================
\begin{document}
\thispagestyle{empty}

% ── Cover block ───────────────────────────────────────────────────────────────
\colorbox{darkgreen}{\parbox{\linewidth}{%
  \vspace{8pt}
  \hspace{8pt}{\color{white}\Large\bfseries 300,000 Streets of Melbourne}\\[3pt]
  \hspace{8pt}{\color{lightgreen}\normalsize School Streets Safety Analysis --- Technical Overview}\\[2pt]
  \hspace{8pt}{\color{white!80!darkgreen}\small Regen Melbourne \texttimes{} RMIT University \quad|\quad
    COSC2667 / COSC2777 \quad|\quad 2024}\\[6pt]
  \hspace{8pt}{\color{white!70!darkgreen}\small
    Siddhartha Ananthula \textbullet{}
    Hishikesh Phukan \textbullet{}
    Reddy Pranay Vipparla \textbullet{}
    Venkata Nagendra Anamala}
  \vspace{8pt}
}}

\vspace{12pt}

% =============================================================================
\section{Problem Statement}

Melbourne's 300,000-plus streets include walking routes used by school children every day. Yet no systematic, data-driven assessment of \emph{which} routes are dangerous exists at scale. Our industry partner, Regen Melbourne, needed a reproducible method to identify high-risk school streets, quantify the severity of hazards, and generate costed intervention recommendations --- all from open government data.

We focused on three secondary schools in the City of Darebin as a proof-of-concept:
\textbf{Preston High School}, \textbf{Reservoir High School}, and \textbf{William Ruthven Secondary College}.

The central research question: \emph{Can data tell us which school streets are putting children at risk --- and does disadvantage make it worse?}

% =============================================================================
\section{Technical Architecture}

The system is a fully reproducible, end-to-end Python pipeline. Every output is regenerated by running a single command: \texttt{python run\_all.py}.

\subsection{Pipeline Stages}

\vspace{4pt}
\begin{tabularx}{\linewidth}{C{0.5cm} L{3.2cm} L{9.5cm}}
\toprule
\rowcolor{darkgreen!90}
\color{white}\textbf{\#} & \color{white}\textbf{Stage} & \color{white}\textbf{What it does} \\
\midrule
\rowcolor{palegreen}
1 & Data Ingestion      & Downloads crash records (VicRoads), street network (OSM via Overpass API), air quality (EPA AirWatch), crime rates (CSA), socio-economic indices (ABS SEIFA 2021), and demographics (ABS Census 2021). \\
2 & Feature Engineering & Computes 26 spatial and environmental features per school gate: footpath coverage, crossing count, shade ratio, OSM amenity density, PM2.5, crime rate per 100k. Sentinel codes (777/888/999) are filtered before any calculation. \\
\rowcolor{palegreen}
3 & HS Scoring          & Ten dedicated scorer modules (\texttt{src/scoring/hs1.py} \ldots \texttt{hs10.py}) each translate raw features into a 0--10 Healthy Streets score using the framework developed by Dr Lucy Saunders (Transport for London, 2015). \\
4 & Severity Classification & \texttt{severity.py} applies rule-based thresholds: \textbf{Major} if HS2~$<$~4.0, HS1~$<$~4.0, or HS5~$<$~3.0; \textbf{Moderate} if two or more indicators fall below 6.0; otherwise \textbf{Minor}. \\
\rowcolor{palegreen}
5 & Recommendations     & A rule engine maps each sub-threshold indicator to a prioritised intervention (Critical / High / Medium / Low) with estimated cost range and delivery timeframe. \\
6 & Machine Learning    & A Ridge regression model (scikit-learn) is trained on all 26 features using Leave-One-Out cross-validation (\textit{n}=3). Identifies which HS scores can be predicted from open data alone. \\
\rowcolor{palegreen}
7 & Scenario Engine     & Ten intervention templates are applied to each school using a delta method: the model estimates the change in HS scores, anchored to actual baseline values. Produces 30 pre-computed before/after results. \\
8 & Equity \& Crash Analysis & SEIFA disadvantage deciles and ABS demographic data are joined to school catchments. Crash records are filtered to 400~m school-gate buffers and broken down by hour. \\
\rowcolor{palegreen}
9 & Visualisation       & Nine static charts (Matplotlib) covering safety scores, equity, crash trends, demographics, feature importance, and ML predictions. GIS layers exported as GeoPackage (\texttt{.gpkg}) and QGIS project. \\
10 & Dashboard Export   & All results, charts, and scenario data are serialised to JSON and embedded in a static GitHub Pages site (Leaflet.js). Zero backend. Zero installation. \\
\bottomrule
\end{tabularx}

\vspace{6pt}

\subsection{Codebase Structure}

\begin{multicols}{2}
\begin{itemize}
  \item \texttt{src/scoring/} --- One module per HS indicator (hs1--hs10) plus severity classifier
  \item \texttt{src/features/} --- \texttt{engineer.py}: all 26 feature computations from OSM, EPA, CSA
  \item \texttt{src/ml/} --- \texttt{train.py}, \texttt{predict.py}, \texttt{evaluate.py}: Ridge model lifecycle
  \item \texttt{src/scenarios/} --- Intervention delta engine; 10 templates
  \item \texttt{src/recommendations/} --- Rule-to-intervention mapping with cost/timeline metadata
  \item \texttt{src/data/} --- Download and cache handlers for each external source
  \item \texttt{src/visualisation/} --- Chart builders; one function per chart
  \item \texttt{docs/} --- Static dashboard (HTML/CSS/JS); deployed to GitHub Pages
  \item \texttt{outputs/} --- All generated files: CSV, PNG, GeoPackage, HTML, PPTX
  \item \texttt{cache/} --- Persistent API response cache; avoids repeated downloads
\end{itemize}
\end{multicols}

% =============================================================================
\section{Analytical Approach}

\subsection{Healthy Streets Framework}

The Healthy Streets framework (Saunders, 2015) provides ten indicators scored 0--10. The minimum ``good'' threshold is 6.0 per indicator. We selected this framework because it is empirically grounded, council-legible, and has been independently validated in multiple cities. Each indicator maps directly to a data source:

\vspace{4pt}
\begin{tabularx}{\linewidth}{C{1.0cm} L{3.5cm} L{4.8cm} L{4.5cm}}
\toprule
\rowcolor{darkgreen!90}
\color{white}\textbf{Code} & \color{white}\textbf{Indicator} & \color{white}\textbf{Measures} & \color{white}\textbf{Data source} \\
\midrule
\rowcolor{palegreen} HS1 & Footpath accessibility & Width, continuity, surface quality & Field survey \\
HS2 & Easy to cross       & Crossing presence, type, visibility & Field survey \\
\rowcolor{palegreen} HS3 & Shade \& shelter       & Tree canopy ratio, shelter density & OpenStreetMap \\
HS4 & Rest places         & Seating, drinking fountain density & OpenStreetMap \\
\rowcolor{palegreen} HS5 & Not too noisy         & Traffic volume, road hierarchy     & Field survey \\
HS6 & Active travel       & Cycling infrastructure, PT proximity & OSM + field \\
\rowcolor{palegreen} HS7 & Feel safe             & Crime rate per 100,000 residents   & CSA Victoria \\
HS8 & Things to see \& do & POI density within 400~m           & OpenStreetMap \\
\rowcolor{palegreen} HS9 & Feel relaxed          & School zone presence, traffic calming & Field survey \\
HS10 & Clean air          & PM2.5 annual average               & EPA AirWatch \\
\bottomrule
\end{tabularx}

\vspace{4pt}

\subsection{Machine Learning Design}

Ridge regression was selected for its robustness to correlated features (OSM spatial features are inherently collinear) and its interpretability for a stakeholder audience. With only three schools, Leave-One-Out cross-validation was the only statistically sound evaluation strategy. Mean MAE across all ten indicators was \textbf{2.88}.

The key ML insight: \textbf{HS7 (feel safe) and HS10 (clean air)} are predictable from open data with low error. \textbf{HS2 (crossing safety)} is not --- because crossing \emph{presence} in OSM does not capture crossing \emph{quality}. This directly informs which indicators require field surveys and which can be automated for a network-wide triage.

\subsection{Equity Analysis}

SEIFA IRSD (Index of Relative Socio-economic Disadvantage) deciles from ABS 2021 were joined to school catchments at SA1 level using population-weighted averaging. The Pearson correlation between SEIFA decile and HS overall score was \textbf{r~=~0.84} across the three schools --- a strong positive relationship suggesting that disadvantaged catchments systematically receive worse walking environments.

% =============================================================================
\section{Key Findings}

\begin{findingbox}
\textbf{Finding 1 --- The Paradox:} Preston High School has the highest overall HS score (7.2/10) but is the \emph{only} Major-severity school. HS2~=~0.4: there is no pedestrian crossing at the school gate on a high-traffic arterial road. Nine green scores cannot compensate for one safety-critical failure.
\end{findingbox}

\begin{insightbox}
\textbf{Finding 2 --- Crashes are rising and school-timed:} Pedestrian and cyclist crashes in Darebin LGA increased 192\% between 2021 and 2024 (25 to 73). At William Ruthven SC, 100\% of crashes within 400~m of the gate occurred during school arrival or pickup windows --- invisible in annual totals, critical when time-filtered.
\end{insightbox}

\begin{insightbox}
\textbf{Finding 3 --- Disadvantage predicts danger (r~=~0.84):} Reservoir High School sits in a Decile~4 catchment (moderate-high disadvantage) and has the most safety deficiencies. Preston High is in Decile~7 and scores highest. The pattern is consistent with Giles-Corti et al.\ (2016): lower-income households have lower car ownership, higher walking dependency, and face worse street conditions.
\end{insightbox}

\begin{insightbox}
\textbf{Finding 4 --- The cheapest fix is a crossing:} A signalised pedestrian crossing at Preston High (cost: \$80k--\$200k, timeframe: $<$1 year) would upgrade HS2 from 0.4 to approximately 3.5, moving severity from Major to Moderate. The scenario engine quantifies this return on investment before councils commit capital.
\end{insightbox}

% =============================================================================
\section{How Findings Were Made Actionable}

Raw scores alone do not move councils. We translated every number into a decision-ready artefact:

\begin{itemize}
  \item \textbf{Severity label} (Major / Moderate / Minor) --- one word that triggers an urgency response without requiring a stakeholder to read ten indicator values.
  \item \textbf{Recommendations table} (14 rows across 3 schools) --- each row includes the specific indicator being addressed, a plain-English description, estimated cost range, and delivery timeframe. Filterable by school and priority in the dashboard.
  \item \textbf{Scenario engine} --- answers the question councils actually ask: ``What do we get for \$80k?'' Thirty pre-computed results are instantly accessible in the browser.
  \item \textbf{Interactive dashboard} (GitHub Pages, zero installation) --- five sections covering hero stats, an interactive 4-layer map (crash spots, walk/cycle networks, KDE heatmap, SEIFA catchments), per-school radar charts, the scenario explorer, and a one-click recommendations CSV download.
  \item \textbf{GIS export} --- \texttt{.gpkg} and \texttt{.qgz} files for council GIS teams to integrate directly into their existing planning workflows.
  \item \textbf{Scalability argument} --- the same pipeline, with no code changes, can be run on all 200+ Melbourne secondary schools. Phase 1 (automated triage on HS7 and HS10) requires zero surveyors.
\end{itemize}

% =============================================================================
\section{Technology Stack}

\begin{tabularx}{\linewidth}{L{3.2cm} L{10.6cm}}
\toprule
\rowcolor{darkgreen!90}
\color{white}\textbf{Layer} & \color{white}\textbf{Tools \& libraries} \\
\midrule
\rowcolor{palegreen} Data \& pipeline & Python 3.12, pandas, geopandas, OSMnx, requests, openpyxl \\
ML & scikit-learn (Ridge, cross\_val\_score), joblib (model persistence) \\
\rowcolor{palegreen} Spatial & Shapely, Fiona, GDAL, GeoPackage, QGIS \\
Visualisation & Matplotlib, Seaborn \\
\rowcolor{palegreen} Dashboard & Leaflet.js, Chart.js, Vanilla JS, GitHub Pages (static) \\
Presentation & python-pptx (automated PPTX generation) \\
\rowcolor{palegreen} Document & \LaTeX{} (book, report, this document) \\
\bottomrule
\end{tabularx}

\vspace{14pt}

\colorbox{darkgreen!90}{\parbox{\linewidth}{%
  \vspace{6pt}
  \hspace{8pt}\color{white}\small
  \textbf{Repository:} github.com/s4111078-sidd/300-000-School-Streets
  \quad\textbullet{}\quad
  \textbf{Dashboard:} s4111078-sidd.github.io/300-000-School-Streets
  \quad\textbullet{}\quad
  \textbf{Run:} \texttt{python run\_all.py}
  \vspace{6pt}
}}

\end{document}
